\section{An actual native norm factor independent of original frames}

\begin{definition}\label{def:native-tc-original-center-open-frame-data}
\lean{EnvelopeProperty.NativeTestConfiguration.TestConfiguration.originalCenterOpenFrame}
\leanok
\uses{thm:native-tc-coherent-all-power-section-norms}
For the original invertible $L$ and actual trivially valued
point $x$, choose one actual original open $U$ containing its
proper center and an actual whole original frame of $L$ on $U$.
This choice is made on the original $X$ before the product
radius, every degree and every original section.
\end{definition}
\begin{proof}
Original test-configuration properness gives properness of
$X$, hence the actual proper center of $x$. The original
invertible module has an original open trivializing cover.
Its union is all of $X$, so some member contains this actual
center. Choose that member and one of its actual whole module
isomorphisms to the original structure sheaf. The resulting
data consist of an original open, an actual frame and the
center-membership proof, without a section evaluation.
\end{proof}

\begin{definition}\label{def:native-tc-intrinsic-native-unit-factor}
\lean{EnvelopeProperty.NativeTestConfiguration.TestConfiguration.gaussNativeNormFactor,EnvelopeProperty.NativeTestConfiguration.TestConfiguration.gaussNativeNormFactor_pos,EnvelopeProperty.NativeTestConfiguration.TestConfiguration.gaussNativeNormFactor_spec}
\leanok
\uses{def:native-tc-original-center-open-frame-data,def:native-tc-coherent-positive-norm-factor,thm:native-tc-coherent-all-power-section-norms}
Define $c_{\mathcal X,L}(x,r)$ by using this actual original
center-open frame in the native transition-unit valuation
construction. It is positive and satisfies the true original
all-degree all-section norm identity. Its definition uses an
actual original native unit and precedes all degrees and
original sections.
\end{definition}
\begin{proof}
The actual completed-point factor through the selected original
open gives membership of the original product Gauss point in
its inverse image. Thus the selected original frame determines
an actual reference residue frame. Take the native unit
transition to the model residue frame induced from the one
proper-center frame of the original total line. Its actual
residue valuation is the defined number. The actual unit
proves positivity, and the coherent original power-norm theorem
applied to these exact original data proves the asserted
identity in every degree and for every original section.
\end{proof}

\begin{theorem}\label{thm:native-tc-intrinsic-factor-original-frame-independence}
\lean{EnvelopeProperty.NativeTestConfiguration.TestConfiguration.gaussNativeNormFactor_eq_compatibleFactor}
\leanok
\uses{def:native-tc-intrinsic-native-unit-factor,cor:native-tc-original-frame-factor-independence,thm:native-tc-unique-positive-all-power-factor}
For the original ample $L$, the defined native factor equals
the native transition valuation constructed using any actual
original center open and whole frame of $L$. Thus its value is
independent of the original center-open frame choice.
\end{theorem}
\begin{proof}
The selected data used in the definition are an actual
original center-open frame. The original frame-independence
theorem compares their native factor with that of any other
actual original center-open frame, giving the asserted
equality. The uniqueness argument uses an actual positive
section in a genuinely produced positive degree, so the
definition and independence require no degree-one section.
\end{proof}
